\section{Introduction}\label{sec:intro}

Cause--effect benchmarks drive method development in causal discovery. For two variables the reference is the T\"ubingen collection \citep{mooij2016distinguishing}, 108 pairs of real measurements of which 102 are univariate; for graphs it is the protein-signalling data of \citet{sachs2005causal} with 17 consensus arcs. A new rule is typically presented in a table in which its number is larger than those of earlier rules. This paper asks how much information such a table carries. It runs a panel of direction rules under one protocol fixed before the runs on eight real pair collections, and asks how large a difference the benchmarks can resolve and why rules that test the conditional structure of a pair \citep{hoyer2009nonlinear,peters2014identifiability,hyvarinen2013pairwise,immer2023loci,lorch2022amortized} so rarely beat simpler rules that read the scale of regression errors or the marginal density \citep{blobaum2018cause,daniusis2010inferring,stienstra2026oneruler}. On T\"ubingen all of them lie between $55\%$ and $67\%$ with intervals of $\pm10$ points; the ordering of the rules on the Causal Chambers \citep{gamella2025chamber} is negatively, but not significantly, correlated with it.

\paragraph{Contributions.} (1) A same-protocol comparison of twelve pair rules on eight real collections (T\"ubingen; six Causal Chambers collections, four of them new; the Sachs pairs), with the stored decision of every rule on every pair, cluster-aware inference and a reproduction of a published headline number (Section~\ref{sec:conventions}). (2) A resolution analysis: a paired test on $102$ pairs has $80\%$ power only for a difference of about $19$ accuracy points, and among ten non-degenerate rules none of $45$ comparisons reaches a raw $p<0.05$; ten reporting universes move a single rule by $11$ to $26$ points and its rank by up to nine places, and $50$ preprocessing configurations by $7$ to $22$ points (Section~\ref{sec:results}). (3) Two unreported implementation conventions, the rounding of a quantiser and the treatment of tied values, that move a rule's accuracy by six and by up to $88$ points (Section~\ref{sec:conventions}). (4) Mechanism evidence, tested by prediction, matched manipulation with a control and leakage-free validation (Section~\ref{sec:mech}): the chamber collections, whose causes are experimenter-set actuators, can be oriented from their marginals alone, a cue that does not transfer to T\"ubingen; additive-noise rules are near-perfect when the noise scale does not depend on the cause and reversed when it does; two attempts to build a rule or a predictive model on this account fail their predictions fixed in advance, which bounds it. (5) Recommendations for reporting (Section~\ref{sec:discussion}), and an accounting of the attempts to build a better rule (Appendix~\ref{app:negative}).

\subsection{Related work}\label{sec:related}
\emph{Direction rules.} Identifiability results for linear non-Gaussian \citep{shimizu2006linear}, additive-noise \citep{hoyer2009nonlinear,peters2014identifiability,zhang2009identifiability} and post-nonlinear models underlie most pair rules; IGCI \citep{daniusis2010inferring} rests on the independence of the cause distribution and the mechanism, RECI \citep{blobaum2018cause} compares regression errors after scaling, LOCI \citep{immer2023loci} fits location-scale models, SLOPE \citep{marx2017slope} uses minimum description length, and learned rules are trained on simulated pairs \citep{lopezpaz2015towards,lorch2022amortized}. \emph{Evaluation.} Real data with known cause--effect structure are scarce: T\"ubingen for pairs, the Causal Chambers \citep{gamella2025chamber} for physical devices with a documented graph, and the Sachs network, a consensus rather than a ground truth. On simulated graphs, \citet{reisach2021beware} showed that the variance structure of additive-noise models can be read by a trivial rule; our marginal-only classifier is a pair-level analogue on real data. A same-hands re-evaluation on T\"ubingen that also introduced the compression baseline we use is \citet{stienstra2026oneruler}; we reproduce its main figure. For graphs we also ran BF-BIC with BOSS \citep{ramsey2025bfbic} and an SNOE-style \citep{huang2025snoe} orientation (Appendix~\ref{sec:graphs}).

\emph{What this paper does not claim.} It does not claim that published methods misreport their results; we audited one published table and reproduced it. It does not claim that dependence-based methods are worse in general: they are among the best rules in the regimes that fit their assumptions and the worst in others, and a benchmark that mixes regimes averages them.
